% GENERATED FILE. Edit canonical lessons, then run npm run generate:books.
% canonical-source-hash: fnv1a64:a24ced82889e624e
% canonical-lessons: HI-C267-rajai, HI-C267-maccharadani, HI-C267-nal, HI-C267-nali, HI-C267-tanki

\chapter{A quilt, A mosquito net, A tap, A drain, A tank}
\label{ch:hi-a-quilt-a-mosquito-net-a-tap-a-drain-a-tank}
\hlchaptermodality{267}

\begin{chapteropening}
\textbf{By the end of this chapter:} \emph{I can say a quilt, a mosquito net, a tap, a drain and a tank.}

Complete the last lesson of chapter 267: I can say a quilt, a mosquito net, a tap, a drain and a tank.
\end{chapteropening}

\section[\texorpdfstring{\hi{रजाई} (rajāī)}{रजाई (rajāī)}]{\hi{रजाई} (rajāī) — a quilt}

\label{lesson:HI-C267-rajai}



\begin{quote}
Before the new one: say the Hindi for a spade, then the Hindi for a mattress.
\end{quote}

\subsection*{You'll want to know: \hi{रजाई}}
\textbf{\hi{रजाई}} — \emph{rajāī} — \textquotedblleft{}a quilt\textquotedblright{}.

\begin{cousinweb}
The origin of \textbf{\hi{रजाई}} (\emph{rajāī}) is uncertain. CDIAL has no entry for it.
\end{cousinweb}

\begin{grammarlens}[title={What you've built}]
One more word for this chapter.
\end{grammarlens}

\subsection*{Your turn}
\begin{itemize}
\raggedright
  \item \emph{Say it:} \emph{rajāī}
  \item \emph{Say it:} \emph{rajāī}, once more
  \item \emph{Say it:} say \emph{rajāī}
  \item \emph{Recall:} say the Hindi for a spade, then the Hindi for a mattress, then say \emph{rajāī} again
\end{itemize}

\begin{tcolorbox}[breakable,colback=teal!4,colframe=teal!35!black,title={Before you move on}]
What does \hi{रजाई} mean? (\textquotedblleft{}A quilt\textquotedblright{}.) Say it once more.
\end{tcolorbox}
\section[\texorpdfstring{\hi{मच्छरदानी} (maccharadānī)}{मच्छरदानी (maccharadānī)}]{\hi{मच्छरदानी} (maccharadānī) — a mosquito net}

\label{lesson:HI-C267-maccharadani}



\begin{quote}
Before the new one: say the Hindi for a mattress, then the Hindi for a quilt.
\end{quote}

\subsection*{You'll want to know: \hi{मच्छरदानी}}
\textbf{\hi{मच्छरदानी}} — \emph{maccharadānī} — \textquotedblleft{}a mosquito net\textquotedblright{}.

\begin{cousinweb}
\textbf{\hi{मच्छरदानी}} (\emph{maccharadānī}) is \emph{macchar}, \textquotedblleft{}mosquito\textquotedblright{}, from Sanskrit \emph{matsara} (CDIAL 9757), with the Persian ending \emph{-dānī}, \textquotedblleft{}holder, container\textquotedblright{}.
\end{cousinweb}

\begin{grammarlens}[title={What you've built}]
One more word for this chapter.
\end{grammarlens}

\subsection*{Your turn}
\begin{itemize}
\raggedright
  \item \emph{Say it:} \emph{maccharadānī}
  \item \emph{Say it:} \emph{maccharadānī}, once more
  \item \emph{Say it:} say \emph{maccharadānī}
  \item \emph{Recall:} say the Hindi for a mattress, then the Hindi for a quilt, then say \emph{maccharadānī} again
\end{itemize}

\begin{tcolorbox}[breakable,colback=teal!4,colframe=teal!35!black,title={Before you move on}]
What does \hi{मच्छरदानी} mean? (\textquotedblleft{}A mosquito net\textquotedblright{}.) Say it once more.
\end{tcolorbox}
\section[\texorpdfstring{\hi{नल} (nal)}{नल (nal)}]{\hi{नल} (nal) — a tap}

\label{lesson:HI-C267-nal}



\begin{quote}
Before the new one: say the Hindi for a quilt, then the Hindi for a mosquito net.
\end{quote}

\subsection*{You'll want to know: \hi{नल}}
\textbf{\hi{नल}} — \emph{nal} — \textquotedblleft{}a tap\textquotedblright{}.

\begin{cousinweb}
\textbf{\hi{नल}} (\emph{nal}) is a tadbhava, a word inherited from Sanskrit and reshaped by sound change. It comes from \emph{naḍa}, \textquotedblleft{}a species of reed\textquotedblright{} (CDIAL 6936). CDIAL lists Hindi \emph{nalā}, \textquotedblleft{}hollow reed, windpipe\textquotedblright{}. It notes possible Dravidian or Munda contact.
\end{cousinweb}

\begin{grammarlens}[title={What you've built}]
One more word for this chapter.
\end{grammarlens}

\subsection*{Your turn}
\begin{itemize}
\raggedright
  \item \emph{Say it:} \emph{nal}
  \item \emph{Say it:} \emph{nal}, once more
  \item \emph{Say it:} say \emph{nal}
  \item \emph{Recall:} say the Hindi for a quilt, then the Hindi for a mosquito net, then say \emph{nal} again
\end{itemize}

\begin{tcolorbox}[breakable,colback=teal!4,colframe=teal!35!black,title={Before you move on}]
What does \hi{नल} mean? (\textquotedblleft{}A tap\textquotedblright{}.) Say it once more.
\end{tcolorbox}
\section[\texorpdfstring{\hi{नाली} (nālī)}{नाली (nālī)}]{\hi{नाली} (nālī) — a drain}

\label{lesson:HI-C267-nali}



\begin{quote}
Before the new one: say the Hindi for a mosquito net, then the Hindi for a tap.
\end{quote}

\subsection*{You'll want to know: \hi{नाली}}
\textbf{\hi{नाली}} — \emph{nālī} — \textquotedblleft{}a drain\textquotedblright{}.

\begin{cousinweb}
\textbf{\hi{नाली}} (\emph{nālī}) belongs with Hindi \emph{nāl}, \textquotedblleft{}hollow stalk, tube\textquotedblright{}, and \emph{nārī}, \textquotedblleft{}tube, vein\textquotedblright{}, which CDIAL lists under Sanskrit \emph{nāḍī}, \textquotedblleft{}tubular stalk, tube\textquotedblright{} (CDIAL 7047).
\end{cousinweb}

\begin{grammarlens}[title={What you've built}]
One more word for this chapter.
\end{grammarlens}

\subsection*{Your turn}
\begin{itemize}
\raggedright
  \item \emph{Say it:} \emph{nālī}
  \item \emph{Say it:} \emph{nālī}, once more
  \item \emph{Say it:} say \emph{nālī}
  \item \emph{Recall:} say the Hindi for a mosquito net, then the Hindi for a tap, then say \emph{nālī} again
\end{itemize}

\begin{tcolorbox}[breakable,colback=teal!4,colframe=teal!35!black,title={Before you move on}]
What does \hi{नाली} mean? (\textquotedblleft{}A drain\textquotedblright{}.) Say it once more.
\end{tcolorbox}
\section[\texorpdfstring{\hi{टंकी} (ṭaṅkī)}{टंकी (ṭaṅkī)}]{\hi{टंकी} (ṭaṅkī) — a tank}

\label{lesson:HI-C267-tanki}



\begin{quote}
Before the new one: say the Hindi for a tap, then the Hindi for a drain.
\end{quote}

\subsection*{You'll want to know: \hi{टंकी}}
\textbf{\hi{टंकी}} — \emph{ṭaṅkī} — \textquotedblleft{}a tank\textquotedblright{}.

\begin{cousinweb}
\textbf{\hi{टंकी}} (\emph{ṭaṅkī}) is a loan from English \emph{tank}. The \emph{-ī} makes a smaller one.
\end{cousinweb}

\begin{grammarlens}[title={What you've built}]
One more word for this chapter.
\end{grammarlens}

\subsection*{Your turn}
\begin{itemize}
\raggedright
  \item \emph{Say it:} \emph{ṭaṅkī}
  \item \emph{Say it:} \emph{ṭaṅkī}, once more
  \item \emph{Say it:} say \emph{ṭaṅkī}
  \item \emph{Recall:} say the Hindi for a tap, then the Hindi for a drain, then say \emph{ṭaṅkī} again
\end{itemize}

\begin{tcolorbox}[breakable,colback=teal!4,colframe=teal!35!black,title={Before you move on}]
What does \hi{टंकी} mean? (\textquotedblleft{}A tank\textquotedblright{}.) Say it once more.
\end{tcolorbox}
